\documentclass[10pt]{article}
\usepackage[utf8]{inputenc}
\usepackage[myheadings]{fullpage}
\DeclareUnicodeCharacter{0301}{\hspace{-1ex}\'{ }}
% Package for headers 
\usepackage{fancyhdr}
\usepackage{lastpage}
\usepackage{listings}

% For figures and stuff
\usepackage{graphicx, wrapfig, subcaption, setspace, booktabs}
\usepackage[T1]{fontenc}
\usepackage{array}

% Change for different font sizes and families
\usepackage[font=small, labelfont=bf]{caption}
\usepackage{fourier}
\usepackage[protrusion=true, expansion=true]{microtype}

% Maths
\usepackage{amsmath,amssymb}
\usepackage{float}
\usepackage{graphicx}
\usepackage{wrapfig}
\usepackage[colorinlistoftodos]{todonotes}
\usepackage[colorlinks=true, allcolors=blue]{hyperref}
\usepackage[super]{nth}

%Grammar expression
\usepackage[altpo]{backnaur}

% Bibliography
\usepackage[
    sorting=none,
	backend=bibtex
]{biblatex} 
\addbibresource{bibliography.bib}

%% Language and font encodings
\usepackage[english]{babel}
\usepackage{csquotes}
\usepackage{titlesec}

\usepackage{appendix}
\renewcommand{\appendixpagename}{Annexes}
\renewcommand{\appendixtocname}{Annexes}
\setcounter{secnumdepth}{4}

\titleformat{\paragraph}
{\normalfont\normalsize\bfseries}{\theparagraph}{1em}{}
\titlespacing*{\paragraph}
{0pt}{3.25ex plus 1ex minus .2ex}{1.5ex plus .2ex}


\newcommand{\HRule}[1]{\rule{\linewidth}{#1}}
\onehalfspacing
\setcounter{tocdepth}{5}
\setcounter{secnumdepth}{5}

%% Sets page size and margins
\usepackage[a4paper,top=2cm,bottom=1.5cm,left=2cm,right=2cm,marginparwidth=1.5cm]{geometry}

\pagestyle{fancy}
\fancyhf{}

% Header and footer information
\setlength\headheight{15pt}
\fancyhead[L]{Thales DIS / EIDD} 
\fancyhead[R]{Djamila Mohamed}
\fancyhead[C]{\today}
\fancyfoot[R]{\thepage}
 \setlength{\marginparwidth}{2cm}


\newcolumntype{P}[1]{>{\centering\arraybackslash}p{#1}}

 
\begin{document}
\begin{singlespace}
\pagenumbering{gobble}

\date{}

% Do not change anything here except in \LARGE \textbf{This is the title of the essay} 
% /hline before and after the title makes those horoziontal lines appear, you can change the appearance by changing the 2pt to different sizees
\title{
		% Change to your faculty if needed
            \includegraphics[scale=0.75]{UniversiteParisCite\_EIDD\_WEB.png}\\   		
            \includegraphics[scale=0.60]{Thales-logo-1.png}\\
            
		\HRule{2pt} \\
		\LARGE \textbf{Implementation of a Frama-C plug-in and a language server for ACSL/C code assistance} %para que quede encerrado en las lineas
		\HRule{2pt} \\ [0.5cm]
		\normalsize \today}
		
\date{}

\author{
          Internship report \\[1cm]
		 Djamila Mohamed      \\
          2023--2024 \\ [1cm]
		 Supervisor:      \\
	   Adel Djoudi, Formal Methods Expert
		 }
		 
\maketitle

\newpage

\tableofcontents
\newpage

\newpage
% Add your own abstract here



\pagenumbering{arabic}
% Uncomment the next line if you want keywords/index terms after the abstract. 
%\textit{\textbf{Keywords}: lorem, ipsum, dolor}
\section{Introduction}
\subsection{Thales DIS, IBS}
Thales Digital Identity and Security is the global business unit (GBU) of Thales which operates in the field of secure software, providing biometric and encryption solutions. Within the DIS GBU is the Identity and Biometric Solutions business line (BL) which is in charge of digital identity (e.g.:\ cards and passports). The IBS BL is divided in three subunits: ID Documents Solutions (IDDS), Digital ID \& Verification Solutions and Public Security. I am assigned to the IDDS subunit as an intern software developer, under the supervision of Adel Djoudi, formal methods expert at Thales DIS.\@

\subsection{Background}
In IDDS, developers use program analyzers such as Frama-C\footnote{https://frama-c.com/html/overview.html}, a C source file analyzer developed by CEA and Inria. This
software provides a precise analysis of their code and potential security flaws in it by expressing functional specifications through ANSI/ISO C Specification Language (ACSL)\footnote{https://frama-c.com/html/acsl.html} annotations. But Frama-C currently does not provide a convenient integrated development environment (IDE) to edit and browse C and ACSL code.

The goal of this project is to contribute to the development of a software solution in OCaml for a high-level support of C/ACSL languages in IDEs including code navigation, completion suggestions, and interaction with Frama-C code analysis.

\section{Research Work}
In this part, we will discuss the main topics covered in the research work I have done so far.
\subsection{The Language Server Protocol}
Developers working on large scale projects often need a code editing environment with the adequate tooling for the languages they use in their workspace. Multiple solutions such as Eclipse or IntelliJ IDEA exist for Java projects, as well as other popular programming languages. But the complexity of developing one tool for one language and one specific editor is tedious as it forces developers to repeat the same logic for every code editor with each programming language. To overcome these, a recent solution was developed: the Language Server Protocol (LSP)\footnote{https://microsoft.github.io/language-server-protocol/overviews/lsp/overview/}. My first task was to do bibliographic research about its background, use cases and some of its implementations to prepare the implementation of the project and understand its challenges.

\subsection{Frama-C plugin development}
We choose to implement the ACSL language server as Frama-C plugin to leverage the parsing features of Frama-C for ACSL code. It was therefore necessary to read and follow the Frama-C (Nickel 28.0) plug-in development guide. The Frama-C plug-in will act as the server for the LSP and the VS Code extension as the client. Most of my tasks revolve around the plug-in development as it is the key part of the project. VS Code was chosen as the default editor since it is used by most developers especially in Thales DIS.\@ As long as we have the server part, any editor can provide an extension to connect to the Frama-C server.

\section{Remaining Tasks}
\begin{itemize}
	\item Initiate tests for the server side
	\item Create the project's proof of concept
	\item Implement the server and add JSON request parsing
	\item Implement VS Code extension
\end{itemize}

\end{singlespace}
\end{document}

